% ---- Back matter: acknowledgements, appendices A--E, references ----
All calculations reported in sect.~9 were performed on a 
dedicated PC cluster at CERN. I am grateful 
to the CERN management for providing the required funds
and to the CERN IT Department for technical support.



\lusapp A. Notational conventions

\vskip-2.5ex

\subsection A.1 Gauge group

The Lie algebra of $\SUn$ 
may be identified with the linear space of all complex, anti-hermitian
traceless $N\times N$ matrices.
With respect to a 
basis $T^a$, $a=1,\ldots,N^2-1$, of such matrices, 
the general element $X$ of the Lie algebra is given by
$X=X^aT^a$ with real components $X^a$
(repeated group indices are automatically summed over).
The generators are assumed to satisfy
\equation{
  \tr\{T^aT^b\}=-\frac{1}{2}\delta^{ab},
  \enum
}
but are otherwise left unspecified.

Gauge fields are normalized such that
gauge transformations and covariant derivatives 
do not involve the gauge coupling.
Lorentz indices $\mu,\nu,\ldots$ in $D$ dimensions 
range from $0$ to $D-1$ and are
summed over when they occur in matching pairs.
The space-time metric is assumed to be Euclidean. In particular,
$p^2=p_{\mu}p_{\mu}$ for any momentum $p$ and $\delta_{\mu\mu}=D$.


\subsection A.2 Dirac matrices

The Dirac matrices $\dirac{\mu}$ satisfy
\equation{
  \diracstar{\mu}{\dagger}=\dirac{\mu},\qquad
  \{\dirac{\mu},\dirac{\nu}\}=2\delta_{\mu\nu}.
  \enum
}
Scalar products $\dirac{\mu}p_{\mu}$
are abbreviated by
$\slash{p}$, as usual, and 
$\sigma_{\mu\nu}={i\over2}\left[\dirac{\mu},\dirac{\nu}\right]$.
In any dimension $D$ close to 4,
\equation{
  \dirac{5}=\dirac{0}\dirac{1}\dirac{2}\dirac{3},
  \qquad
  \tr\{1\}=4.
  \enum
}
Note that the axial currents require renormalization if dimensional
regularization is used with these conventions for the Dirac matrices.


\subsection A.3 Lattice theory

The lattice theories considered in this paper are 
set up as usual on a four-dimen\-sional hypercubic 
lattice with spacing $a=1$. 
For the quark fields, Wilson's formulation is chosen,
where the quark fields reside on the sites of the lattice
and carry the same indices as the quark fields in the 
continuum theory.
The gauge covariant forward and backward difference operators
acting on a quark field $\psi(x)$ in presence of 
the gauge field $U(x,\mu)$ are defined by
\equation{
  \nab{\mu}\psi(x)=U(x,\mu)\psi(x+\hat{\mu})-\psi(x),
  \enum
  \nexteq{2.5ex}
  \nabstar{\mu}\psi(x)=\psi(x)-U(x-\hat{\mu},\mu)^{-1}\psi(x-\hat{\mu}),
  \enum
}
where $\hat{\mu}$ denotes the unit vector in direction $\mu$.
On the lattice, $\drv{\mu}$ and $\drvstar{\mu}$ stand for
the ordinary forward and backward difference operators 
and
\equation{
  \ring{\partial}_{\mu}=\frac{1}{2}(\drv{\mu}+\drvstar{\mu})
  \enum
}
for their symmetric combination.

The differential operators $\partial^a_{x,\mu}$ act on 
differentiable functions $f(U)$ of the gauge field $U$ according to
\equation{
  \partial^a_{x,\mu}f(U)=
  {\rmd\over\rmd s}f(\rme^{sX}U)\!\left.{{\vphantom{r\over s}}}\right|_{s=0},
  \hskip1.7em
  X(y,\nu)=\cases{T^a & if $(y,\nu)=(x,\mu)$,\cr
                          \noalign{\vskip0.8ex}
                        0  & otherwise.\cr}
  \enum
}
While these operators depend on the choice of the generators $T^a$,
the gradient field
\equation{
  (\partial f)(U,x,\mu)=T^a\partial^a_{x,\mu}f(U)
  \enum
}
can be shown to be basis-independent.



\lusapp B. Properties of the mapping (5.9)

The mapping (5.9) is differentiable and satisfies
\equation{
  \Psu(\rme^{sX})\mathrel{\mathop=_{s\to0}}
  sX+\rmO(s^2)
  \enum
}
for all matrices $X$ in the Lie algebra of $\SUn$. 
These properties imply that
there exists an open neighbourhood of unity in $\SUn$,
where $\Psu$ is one-to-one and in which
the Jacobian matrix
\equation{
  J(U)^{ab}=-2\kern1pt\tr\kern-1pt\left\{T^a\Psu(T^bU)\right\}
  \enum
} 
is nowhere singular.

Using simple norm estimates, 
the interior of the domain
\equation{
  \Dsu=\left\{U\in\SUn\bigm|
  \Re\tr\{1-U\}\leq(16N)^{-1}\right\}
  \enum
}
can be shown to be such a neighbourhood of unity. Moreover,
the Jacobian $\det J(U)$ is strictly positive 
on $\Dsu$.
As a consequence, the identity
\equation{
  \int_{\Dsu} \rmd U\,\delta\left(\Psu(U)-\Psu(V)\right)f(U)=
  k\det J(V)^{-1}f(V)
  \enum
}
holds for all integrable functions $f(U)$
and all $V\in\Dsu$, the Dirac $\delta$-function being the 
one associated to the $\SUn$ invariant measure 
on the Lie algebra.
The proportionality 
constant $k$ in this equation is positive and could be worked out
explicitly, but its value is not needed in this paper.



\lusapp C. Wick contractions

The Wick contractions of the basic fermion fields 
in the O($a$) improved lattice theory 
are first worked out for the case of the 
discretized flow equations with time step $\eps$.
Time-ordering ambiguities are avoided in this way and
the limit $\eps\to0$ can then easily be taken at the end of 
the calculation.

\subsection C.1 Fermion action and functional integral

To get started, it may be helpful to recall that
the fermion part of the total action of the O($a$) improved lattice 
theory in 4+1 dimensions is the sum of
a boundary term,
\equation{
  \SF=\sum_x\bigl\{\psibar(x)\Dm\psi(x)
  +\cfl\lambdabar(0,x)\lambda(0,x)\bigr\},
  \enum
}
and of the bulk action (5.8).
The Wilson--Dirac operator $\Dm$ in eq.~(C.1) includes the
Sheikholeslami--Wohlert term [\cite{SW},\cite{SFimp}] and the 
mass term with bare mass matrix $M_0$.

In the bulk action (5.8), it is understood that the 
time-dependent fields satisfy the boundary conditions (2.6).
The unconstrained fermion (and similarly the anti-fermion) fields 
integrated over in the functional integral can thus be
taken to be 
\equation{
  \psi(x),\quad\chi(\eps,x),\ldots,\chi(T,x),\quad
  \lambda(0,x),\ldots,\lambda(T-\eps,x).
  \enum
}
Since the action is quadratic in these fields,
the fermion integral is given through Wick's theorem.


\subsection C.2 Field transformation

In the following, an important r\^ole is played by
the kernel $K_{\eps}(t,x;s,y)$ that satisfies
\equation{
  K_{\eps}(t,x;t,y)=\delta_{xy}
  \enum
}
at all $t$ in the range $0\leq t\leq T$ and 
the discretized quark flow equation
\equation{
  \left(\partial_t-\Delta\right)K_{\eps}(t,x;s,y)=0
  \enum
}
if $0\leq s\leq t<T$.
Clearly, since the flow equation (C.4) can be solved 
step by step starting from time $t=s$,
these equations uniquely determine the kernel at
all $t\geq s$.

The calculation of the fermion propagators can now be simplified by
performing a field transformation, where 
$\chi(t,x)$ is expressed through $\psi(x)$ and 
another field $\phi(t,x)$, defined at $t=\eps,\ldots,T$,
according to
\equation{
  \chi(t,x)=\sum_yK_{\eps}(t,x;0,y)\psi(y)
  +\eps\sum_{s=\eps}^t\sum_yK_{\eps}(t,x;s,y)\phi(s,y).
  \enum
}
For any given field $\psi(x)$,
the mapping from the field variables
$\chi(\eps,x),\ldots,\chi(T,x)$ to the new variables 
$\phi(\eps,x),\ldots,\phi(T,x)$ is invertible in view of the 
identity
\equation{
  \phi(t+\eps,x)=\cases{
  \left(\partial_t-\Delta\right)\chi(t,x) & if $0<t<T$,\cr
  \noalign{\vskip1.5ex}
  {1\over\eps}(\chi(\eps,x)-\psi(x))-\Delta\psi(x) &
  if $t=0$.\cr
  }
  \enum
}
Moreover, the associated Jacobian matrix is lower triangular
and its determinant does not depend on the gauge field.

After the field transformation and the corresponding
transformation of the anti-fermion fields,
the fermion action assumes the form
\equation{
  \SF+\SFfl=\sum_x
  \bigl\{\psibar(x)\Dm\psi(x)+\cfl\lambdabar(0,x)\lambda(0,x)\bigr\}
  \noenum
  \nexteq{0.0ex}
  {\phantom{\SF+\SFfl=}}
  +\eps\sum_{t=0}^{T-\eps}\sum_x\bigl\{
  \lambdabar(t,x)\phi(t+\eps,x)+\phibar(t+\eps,x)\lambda(t,x)\bigr\}.
  \enum
}
The transformation thus achieves a nearly perfect decoupling of the 
field variables.


\subsection C.3 Contractions

Since the fundamental quark fields $\psi,\psibar$ are completely
decoupled from the other fields, the contraction of these fields 
\equation{
  \wick{\psi(x)\psibar(y)}=S(x,y),
  \enum
  \nexteq{2.5ex}
  \Dm S(x,y)=\delta_{xy},
  \enum
}
coincides with the quark propagator in the theory in 4 dimensions
as expected.
The other non-vanishing contractions deriving 
from the action (C.7) are
\equation{
  \verylongwick{\phi(t+\eps,x)\lambdabar(t,y)}=
  \longwick{\lambda(t,x)\phibar(t+\eps,y)}={1\over\eps}\delta_{xy},
  \quad t=0,\eps,\ldots T-\eps,
  \enum
  \nexteq{2.5ex}
  \longwick{\phi(\eps,x)\phibar(\eps,y)}=-{\cfl\over\eps^2}\delta_{xy}.
  \enum
}
For the contractions involving the fields $\chi$ and $\chibar$
one then obtains
\equation{
  \longwick{\chi(t,x)\lambdabar(s,y)}=
  \theta(t>s)K_{\eps}(t,x;s+\eps,y),
  \enum
  \nexteq{3.0ex}
  \longwick{\lambda(t,x)\chibar(s,y)}=
  \theta(t<s)K_{\eps}(s,y;t+\eps,x)^{\dagger},
  \enum
  \nexteq{3.0ex}
  \longwick{\chi(t,x)\psibar(y)}=\sum_{v}
  K_{\eps}(t,x;0,v)S(v,y),
  \enum
  \nexteq{2.4ex}
  \wick{\psi(x)\chibar(s,y)}=\sum_{w}
  S(x,w)K_{\eps}(s,y;0,w)^{\dagger},
  \enum
  \nexteq{2.4ex}
  \longwick{\chi(t,x)\chibar(s,y)}=\sum_{v,w}
  K_{\eps}(t,x;0,v)S(v,w)
  K_{\eps}(s,y;0,w)^{\dagger}
  \noenum
  \nexteq{1.5ex}
  {\phantom{\longwick{\chi(t,x)\chibar(s,y)}=}}
  -\cfl\sum_vK_{\eps}(t,x;\eps,v)K_{\eps}(s,y;\eps,v)^{\dagger}.
  \enum  
}
In all these equations, it is understood that the flow-time arguments of 
$\chi$ and $\chibar$ are in the range $(0,T]$ and those of 
$\lambda$ and $\lambdabar$ in the range $[0,T)$. 


\subsection C.4 Contractions in the continuous time limit

When the time step $\eps$ is taken to zero, the kernel $K_{\eps}(t,x;s,y)$
smoothly converges to the fundamental solution $K(t,x;s,y)$
of the quark flow equation in the continuous time formulation of
the lattice theory.
The complete set of the non-zero contractions of the fermion
fields, given by eqs.~(C.8) and (C.12)--(C.16), therefore has a 
well-defined limit.










\lusapp D. Integration of the lattice flow equations

The numerical integration of the pure-gauge gradient flow has already been 
discussed in appendix C of ref.~[\cite{WilsonFlow}]. 
Here the 3rd order Runge--Kutta integrator described there 
is extended to the quark flow equation.


\subsection D.1 Integration of the gradient flow

Following ref.~[\cite{WilsonFlow}],
the flow equation (5.2) is written
in an abstract form
\equation{
   \partial_tV_t=Z(V_t)V_t,
   \enum
}
where the gauge field $V_t$
is considered to be an element of a 
(high-dimensional) Lie group and the gradient $Z(V_t)$
of the Wilson action an element of the associated Lie algebra.
The Runge--Kutta integrator mentioned above 
proceeds in time steps of size $\eps$.
Assuming $V_t$ is known at some flow time $t$, an approximation to 
the exact solution $V_{t+\eps}$ at time $t+\eps$ 
is obtained by computing the fields
\equation{
  W_0=V_t,
  \noenum
  \nexteq{2.5ex}
  W_1=\exp\bigl\{\frac{1}{4}Z_0\bigr\}W_0,
  \noenum
  \nexteq{2.5ex}
  W_2=\exp\bigl\{\frac{8}{9}Z_1-\frac{17}{36}Z_0\bigr\}W_1,
  \noenum
  \nexteq{2.5ex}
  W_3=\exp\bigl\{\frac{3}{4}Z_2
                 -\frac{8}{9}Z_1+\frac{17}{36}Z_0\bigr\}W_2,
  \enum
}
where
\equation{
  Z_i=\eps Z(W_i),\qquad i=0,1,2.
  \enum
}
These rules are such that
\equation{
  W_3=V_{t+\eps}+\rmO(\eps^4)
  \enum
}
and $W_3$ thus provides the desired approximation to $V_{t+\eps}$.


\subsection D.2 Integration of the quark flow equation

It is again helpful to rewrite
the flow equation in an abstract form,
\equation{
  \partial_t\chi_t=\Delta(V_t)\chi_t,
  \enum
}
where the quark field $\chi_t$
is considered to be an element of a 
complex vector space. 
Given $V_t$ and $\chi_t$ at some flow time $t$, and assuming
the gauge fields $W_0$, $W_1$ and $W_2$ are as above,
the quark fields
\equation{
  \phi_0=\chi_t,
  \noenum
  \nexteq{2.5ex}
  \phi_1=\phi_0+\frac{1}{4}\Delta_0\phi_0,
  \noenum
  \nexteq{2.5ex}
  \phi_2=\phi_0+\frac{8}{9}\Delta_1\phi_1-\frac{2}{9}\Delta_0\phi_0,
  \noenum
  \nexteq{2.5ex}
  \phi_3=\phi_1+\frac{3}{4}\Delta_2\phi_2,
  \enum
}
may then be calculated, where
\equation{
  \Delta_i=\eps\Delta(W_i),\qquad i=0,1,2.
  \enum
}
It is not difficult to show that
\equation{
  \phi_3=\chi_{t+\eps}+\rmO(\eps^4),
  \enum
}
and $\phi_3$ thus approximates $\chi_{t+\eps}$ 
to an accuracy that matches the one attained by 
the integrator of the pure-gauge flow.


\lusapp E. Integration of the adjoint flow equation (7.10)

The numerical integration of the adjoint flow equation
is complicated by the fact that the 
backward integration of the flow equation for the gauge field
is exponentially unstable. 
In this appendix, the problem is first reformulated
and its solution is then discussed. For simplicity, the index
$k$ in eq.~(7.10) is omitted.


\subsection E.1 Adjoint Runge--Kutta integration

For a given step size
$\eps>0$ and $t=n\eps$, $n=0,1,2,\ldots$, let
\equation{
  V_t^{\eps}(x,\mu)\quad\hbox{and}\quad
  \chi_t^{\eps}(x)
  \enum
}
be the gauge and quark fields generated by the Runge--Kutta integrator
defined in appendix D. There exists another sequence of quark fields
\equation{
  \xi_s^{\eps}(x),\quad s=t,t-\eps,t-2\eps,\ldots,0,
  \enum
}
with initial value
\equation{
  \xi_t^{\eps}(x)=\eta(x),
  \enum
}
such that the identity
\equation{
  (\xi_s^{\eps},\chi_s^{\eps})=(\eta,\chi_t^{\eps})
  \enum
}
holds exactly for all $s$. This property ensures that
$\xi_s^{\eps}$ is an approximate solution of the adjoint
flow equation, accurate to order $\eps^3$, 
with the required initial value.

The Runge--Kutta steps that lead from $\chi_s^{\eps}$
to $\chi_{s+\eps}^{\eps}$ amount to applying the operator
\equation{
  R_s^{\eps}=1+\frac{1}{4}\Delta_0+\frac{3}{4}\Delta_2
  \bigl\{1-\frac{2}{9}\Delta_0+\frac{8}{9}\Delta_1(1+\frac{1}{4}\Delta_0)
  \bigr\}
  \enum
}
to $\chi_s^{\eps}$, where $\Delta_i$ is given by eq.~(D.7) 
and the fields $W_i$ that appear there 
by eq.~(D.2) with $W_0=V_s^{\eps}$.
From eq.~(E.4) one then infers that
\equation{
  \xi_s^{\eps}(x)=({R_s^{\eps}}^{\dagger}\xi_{s+\eps}^{\eps})(x)
  \enum
  \nexteq{2.5ex}
  {R_s^{\eps}}^{\dagger}=1+\frac{1}{4}\Delta_0+
  \bigl\{1-\frac{2}{9}\Delta_0+(1+\frac{1}{4}\Delta_0)\frac{8}{9}\Delta_1
  \bigr\}\frac{3}{4}\Delta_2,
  \enum
}
and $\xi_s^{\eps}$ is thus obtained from 
$\xi_{s+\eps}^{\eps}$ by calculating
\equation{
  \lambda_3=\xi_{s+\eps}^{\eps},
  \noenum
  \nexteq{2.5ex}
  \lambda_2=\frac{3}{4}\Delta_2\lambda_3,
  \noenum
  \nexteq{2.5ex}
  \lambda_1=\lambda_3+\frac{8}{9}\Delta_1\lambda_2,
  \noenum
  \nexteq{2.5ex}
  \lambda_0=\lambda_1+\lambda_2+
  \frac{1}{4}\Delta_0(\lambda_1-\frac{8}{9}\lambda_2),
  \enum
}
and setting $\xi_s^{\eps}=\lambda_0$.


\subsection E.2 Stability issue

The adjoint recursion (E.8) is a smoothing operation and is therefore
numerically stable. To be able to apply the recursion at time $s+\eps$,
the gauge field $V_s^{\eps}$ at flow time $s$ must however be known.
In principle, the sequence of gauge fields from $s=t$ to $s=0$
could be reconstructed through
backward integration of the gradient flow, 
but such strategies are not practical,
because the flow is exponentially unstable in this direction.

Alternatively, $V_s^{\eps}$ can be computed through
forward integration of the flow starting from the fundamental
field $V_0^{\eps}=U$. If the field is recomputed for every $s$,
the total number of gauge update steps required in the 
course of the integration
scales like the square of the number $m=t/\eps$ of time steps.
This method is numerically safe and the calculated
quark fields satisfy the identity (E.4) practically to machine precision.
The required computational effort can however be prohibitively large.


\subsection E.3 Improved scheme

The adjoint flow equation can be integrated more efficiently
by dividing the sequence $0,\eps,\ldots,t$ of flow
times into $n_b$ consecutive blocks of about equal size.
One may then first compute the field $V_r^{\eps}$ at 
the time $r$ at the beginning of the last block 
by forward integration from time 0 and
keep this field in the memory of the computer. After that
the adjoint flow equation is integrated from time $t$ to $r$
by applying the update step (E.8) and 
using the safe reconstruction method described above
for the gauge field,
where the reconstruction now starts 
from the stored field rather than the field
at time $0$.

In the next step the gauge field is calculated
at the first time $r$ in the next-to-last block 
and stored in memory. The integration of the 
adjoint flow equation can then proceed to time $r$
as before. Clearly, the blocks can 
be processed in this way one after another
until the integration reaches time $0$.

The total number of gauge update steps
in this scheme is minimized if $m=n_b^2$.
While $m$ will rarely be the square of an integer,
one can always set $n_b=\lceil m^{1/2}\rceil$ and divide
$m$ into blocks of size $m_b=\lfloor m/n_b\rfloor$ and $m_b+1$.
The computational effort then scales approximately like
$m^{3/2}$ and is much smaller than the one required
for the safe scheme described in the previous subsection
already at low values of $m$.


\subsection E.4 Hierarchical scheme

A further acceleration of the computation can be achieved 
by dividing each block
into $n_b$ smaller blocks, and these into $n_b$ even smaller
blocks, and so on, with a
saved gauge field at each block level.
Given $m$ and the number $n_s$ of saved gauge fields, 
a nearly optimal choice of the block number $n_b$ is
\equation{
  n_b=\lceil m^{1/(n_s+1)}\rceil.
  \enum
}
For $m\leq10^3$ and $n_s=4$, for example, the
total number of gauge-field updates that need to be 
performed is then not more than about $5m$.
Moreover, 
when several quark fields are evolved simultaneously,
the intermediate gauge fields need to be computed only once.

\bigskip
\noindent{\large\bfseries References}
\medskip
\begin{thebibliography}{99}



% Wilson flow & emergence of the topological sectors

\bibitem{WilsonFlow}
M. L\"uscher,
{\it Properties and uses of the Wilson flow in lattice QCD},
JHEP 1008 (2010) 071


% Renormalization of the gradient flow

\bibitem{RenFlow}
M. L\"uscher, P. Weisz,
{\it Perturbative analysis of the gradient flow in non-Abelian gauge
theories},
JHEP 1102 (2011) 051


% Step scaling

\bibitem{StepScaling}
M. L\"uscher, P. Weisz, U. Wolff,
{\it A numerical method to compute the running coupling in asymptotically 
free theories},
Nucl. Phys. B359 (1991) 221

\bibitem{NogradiEtAl}
Z. Fodor, K. Holland, J. Kuti, D. Nogradi, C. H. Wong,
{\it The Yang-Mills gradient flow in finite volume},
JHEP 1211 (2012) 007; {\it The gradient flow running coupling scheme},
PoS (Lattice 2012) 050

\bibitem{FritzschRamos}
P. Fritzsch, A. Ramos,
{\it The gradient flow coupling in the Schr\"odinger functional},
arXiv:1301.4388


% Quark source smearing

\bibitem{Guesken}
S. G\"usken,
{\it A study of smearing techniques for hadron correlation functions},
Nucl. Phys. B (Proc. Suppl.) 17 (1990) 361

\bibitem{AlexandrouEtAl}
C. Alexandrou, F. Jegerlehner, S. G\"usken, K. Schilling, R. Sommer,
{\it B meson properties from lattice QCD},
Phys. Lett. B256 (1991) 60


% Renormalization of the Langevin equation

\bibitem{Zinn}
J. Zinn--Justin,
{\it Renormalization and stochastic quantization},
Nucl. Phys. B275 [FS17] (1986) 135

\bibitem{ZinnZwanziger}
J. Zinn--Justin, D. Zwanziger, 
{\it Ward identities for the stochastic quantization of gauge fields},
Nucl. Phys. B295 [FS21] (1988) 297


% Renormalization of quantum field theories on a half space

\bibitem{Symanzik}
K. Symanzik,
{\it Schr\"odinger representation and Casimir effect in
renormalizable quantum field theory},
Nucl. Phys. B190 [FS3] (1981) 1


% Wilson formulation of lattice QCD

\bibitem{Wilson}
K. G. Wilson, {\it Confinement of quarks}, Phys. Rev. D10 (1974) 2445


% Exact chiral symmetry on the lattice

% Ginsparg-Wilson relation

\bibitem{GinspargWilson}
P. H. Ginsparg, K. G. Wilson,
{\it A remnant of chiral symmetry on the lattice},
Phys. Rev. D25 (1982) 2649


% Domain wall fermions

\bibitem{Kaplan}
D. B. Kaplan,
{\it A method for simulating chiral fermions on the lattice},
Phys. Lett. B288 (1992) 342;
{\it Chiral fermions on the lattice},
Nucl. Phys. B (Proc. Suppl.) 30 (1993) 597

\bibitem{Shamir}
Y. Shamir,
{\it Chiral fermions from lattice boundaries},
Nucl. Phys. B406 (1993) 90

\bibitem{FurmanShamir}
V. Furman, Y. Shamir,
{\it Axial symmetries in lattice QCD with Kaplan fermions},
Nucl. Phys. B439 (1995) 54


% Classically perfect Dirac operator

\bibitem{Hasenfratz}
P. Hasenfratz, {\it Prospects for perfect actions},
Nucl. Phys. B (Proc. Suppl.) 63 (1998) 53;
{\it Lattice QCD without tuning, mixing and current renormalization},
Nucl. Phys. B525 (1998) 401

\bibitem{HLN}
P. Hasenfratz, V. Laliena, F. Niedermayer,
{\it The index theorem in QCD with a finite cutoff},
Phys. Lett. B427 (1998) 125


% Neuberger-Dirac operator

\bibitem{NeubergerDirac}
H. Neuberger,
{\it Exactly massless quarks on the lattice},
Phys. Lett. B417 (1998) 141;
{\it More about exactly massless quarks on the lattice},
{\it ibid.}\/ B427 (1998) 353


% Lattice chiral symmetry

\bibitem{ExactChSy}
M. L\"uscher,
{\it Exact chiral symmetry on the lattice and the Ginsparg--Wilson relation},
Phys. Lett. B428 (1998) 342

\bibitem{BoulderReview}
F. Niedermayer,
{\it Exact chiral symmetry, topological charge and related topics},
Nucl. Phys. B (Proc. Suppl.) 73 (1999) 105


% SU(2) chiral perturbation theory

\bibitem{ChPT}
J. Gasser, H. Leutwyler,
{\it Chiral perturbation theory to one loop},
Ann. Phys. 158 (1984) 142


% Global existence of the smoothing flow

\bibitem{ArnoldI}
V. I. Arnold,
{\it Ordinary Differential Equations}, 3rd ed. (Springer-Verlag, Berlin, 2008)


% Stout smearing

\bibitem{Stout}
C. Morningstar, M. Peardon, 
{\it Analytic smearing of SU(3) link variables in lattice QCD},
Phys. Rev. D69 (2004) 054501


% Symanzik effective action

\bibitem{SymanzikSeffI}
K. Symanzik, 
{\it Some topics in quantum field theory},
in: Mathematical problems in theoretical physics, eds. R. Schrader
et al., Lecture Notes in Physics, Vol. 153 (Springer, New York, 1982)

\bibitem{SymanzikSeffII}
K. Symanzik,
{\it Continuum limit and improved action in lattice theories:
(I). Principles and $\varphi^4$ theory},
Nucl. Phys. B226 (1983) 187


% On-shell improvement

\bibitem{OnShell}
M. L\"uscher, P. Weisz,
{\it On-shell improved lattice gauge theories},
Commun. Math. Phys. 97 (1985) 59 [E: {\it ibid.} 98 (1985) 433]


% O(a) improvement

\bibitem{SW}
B. Sheikholeslami, R. Wohlert, 
{\it Improved continuum limit lattice action for QCD with Wilson fermions},
Nucl. Phys. B259 (1985) 572

\bibitem{SFimp}
M. L\"uscher, S. Sint, R. Sommer, P. Weisz,
{\it Chiral symmetry and O(a) improvement in lattice QCD},
Nucl. Phys. B478 (1996) 365


% O(a) improvement with non-degenerate quarks

\bibitem{ImpNonD}
T. Bhattacharya, R. Gupta, W. Lee, S. R. Sharpe, J. M. S. Wu,
{\it Improved bilinears in lattice QCD with nondegenerate quarks},
Phys. Rev. D73 (2006) 034504


% Random sources

\bibitem{MichaelPeisa}
C. Michael, J. Peisa,
{\it Maximal variance reduction for stochastic propagators with 
applications to the static quark spectrum},
Phys. Rev. D58 (1998) 034506


% QCD with open boundary conditions and tm reweighting

\bibitem{TMopenQCD}
M. L\"uscher, S. Schaefer,
{\it  	
Lattice QCD with open boundary conditions and twisted-mass reweighting},
Comput. Phys. Commun. 184 (2013) 519


% Open boundary conditions

\bibitem{openQCD}
M. L\"uscher, S. Schaefer,
{\it Lattice QCD without topology barriers},
JHEP 1107 (2011) 036


% PACS-CS physical point simulations

\bibitem{AokiEtAlI}
S. Aoki et al.~(PACS-CS collab.), 
{\it 2+1 flavor lattice QCD toward the physical point},
Phys. Rev. D79 (2009) 034503

\bibitem{AokiEtAlII}
S. Aoki et al.~(PACS-CS collab.),
{\it Physical point simulation in 2+1 flavor lattice QCD},
Phys. Rev. D81 (2010) 074503


% Iwasaki action

\bibitem{Iwasaki}
Y. Iwasaki,
{\it Renormalization group analysis of lattice theories and
improved lattice action. II -- Four-dimensional non-Abelian SU(N) gauge model},
preprint UTHEP-118 (1983) and arXiv:1111.7054


% Non-perturbative O(a) improvement

\bibitem{AokiEtAlIII}
S. Aoki et al. (CP-PACS collab.),
{\it Nonperturbative O(a) improvement of the Wilson quark action with the 
RG-improved gauge action using the Schr\"odinger functional method},
Phys. Rev. D73 (2006) 034501

\bibitem{KanekoEtAl}
T. Kaneko et al. (CP-PACS/JLQCD and ALPHA collab.),
{\it Non-perturbative improvement of the axial current with 
three dynamical flavors and the Iwasaki gauge action},
JHEP 0704 (2007) 092


% Computation of ZA and ZP

\bibitem{AokiEtAlIV}
S. Aoki et al. (PACS-CS collab.),
{\it Non-perturbative renormalization of quark mass 
in $N_f$=2+1 QCD with the Schr\"odinger functional scheme},
JHEP 1008 (2010) 101


% SF method for ZA

\bibitem{ZASFI}
M. L\"uscher, S. Sint, R. Sommer, H. Wittig,
{\it Non-perturbative determination of the axial current 
normalization constant in O(a) improved lattice QCD},
Nucl. Phys. B491 (1997) 344

\bibitem{ZASFII}
M. Della Morte, R. Hoffmann, F. Knechtli, R. Sommer, U. Wolff,
{\it Non-pertur\-ba\-tive renormalization of the axial current 
with dynamical Wilson fermions},
JHEP 0507 (2005) 007





\end{thebibliography}
